% Figure 53.1 : la politique de décalage de versions, pour un API server en 1.37 (chapitre 53).
\documentclass[tikz,border=4pt]{standalone}
\usepackage{quai}
\begin{document}
\begin{tikzpicture}[x=1.6cm,y=1cm,
    nom/.style={font=\scriptsize, anchor=east, align=right, text width=4.2cm},
    barre/.style 2 args={draw=#1, fill=#2, rounded corners=1.5pt, minimum height=0.5cm},
    lien/.style={qlbl, font=\scriptsize}]
  \foreach \v [count=\i from 0] in {1.34,1.35,1.36,1.37,1.38} {
    \node[font=\scriptsize\ttfamily, text=QMuted] at (\i,0.6) {\v};
    \draw[draw=QGrisBg, line width=0.6pt] (\i,0.35) -- (\i,-4.9);
  }
  \foreach \y/\n/\a/\b/\c/\note in {
    0/{\textbf{kube-apiserver}\\la référence}/3/3/QOcre/{N},
    -1/{\textbf{kube-controller-manager},\\\textbf{kube-scheduler}, \textbf{cloud-controller-manager}}/2/3/QViolet/{N ou N-1},
    -2/{\textbf{kubelet}}/0/3/QBleu/{N à N-3},
    -3/{\textbf{kube-proxy}}/0/3/QSarcelle/{N à N-3 (et trois versions d'écart au plus avec son kubelet)},
    -4/{\textbf{kubectl}}/2/4/QVert/{N-1 à N+1}} {
    \node[nom] at (-0.35,\y) {\n};
    \node[barre={\c}{\c Bg}, minimum width={(\b-\a)*1.6cm+0.5cm}, anchor=west] at ({\a-0.16},\y) {};
    \node[lien, anchor=west] at (4.35,\y) {\note};
  }
  \node[qlbl, font=\scriptsize, anchor=north west, text width=12cm, align=left] at (-2.9,-5.1) {Avec plusieurs API servers (haute disponibilité), les plus récents et les plus anciens ne peuvent différer que d'une version mineure : c'est pourquoi on monte d'une version mineure à la fois.};
\end{tikzpicture}
\end{document}
